% twin-axes: temperature and pressure of one batch against time, in the classic (Matplotlib) style, polished:
% ax.twinx() with each y axis in the colour of its series, the three phases of the run as axvspan bands
% labelled at the top (heating, isothermal hold, cooling) and the set point as a dashed axhline. No legend: the
% coloured axes are the key. The pressure axis runs to 6 bar so its plateau sits below the temperature plateau
% instead of on it. The controller's overshoot and the decaying oscillation after the ramp are in the data
% (data/twin-axes.dat, from data/make_data.py).
\documentclass[10pt,tikz,border=4pt]{standalone}
\usepackage{classicfig}
\begin{document}
\begin{tikzpicture}
\begin{axis}[classic, twin color=C0, mpl grid light, mpl minor={1}{1}, mpl title left,
  title={Reactor log}, xlabel={Time (h)}, ylabel={Temperature (\textdegree C)},
  xmin=-0.6, xmax=24.6, enlarge x limits=false, ymin=15, ymax=85, enlarge y limits=false,
  xtick={0,4,...,24}, ytick={20,30,...,80}, ytick pos=left]
\mpltitleright{batch R-17, 2.5 L autoclave}
% phases of the run, behind everything
\axvspan[C1, opacity=0.09]{0}{3}
\axvspan[C0, opacity=0.07]{18}{24}
\node[mpl small, text=black!60, anchor=north] at (axis cs:1.5,83) {heating};
\node[mpl small, text=black!60, anchor=north] at (axis cs:10.5,83) {isothermal hold};
\node[mpl small, text=black!60, anchor=north] at (axis cs:21,83) {cooling};
% set point
\axhline[C0!70, line width=0.8pt, dash pattern=on 4pt off 2.5pt]{70}
\node[mpl small, text=C0!85!black, anchor=north east] at (axis cs:17.7,69.4) {set point 70 \textdegree C};
\addplot[C0] table[x=t, y=T] {data/twin-axes.dat};
\end{axis}
\begin{axis}[twin right, twin color=C3,
  ylabel={Pressure (bar)},
  xmin=-0.6, xmax=24.6, enlarge x limits=false, ymin=0, ymax=6, enlarge y limits=false, ytick={0,...,6},
  minor y tick num=1]
\addplot[C3] table[x=t, y=P] {data/twin-axes.dat};
\end{axis}
\end{tikzpicture}
\end{document}
